\chapter{TandemApp: the dashboard and proxy on their own}\label{ch:install-app}

The dashboard\index{dashboard} also runs without the engine, as TandemApp\index{TandemApp} (\code{tandem-app}), in front of any OpenAI-compatible engine or cloud API. Clients send their requests to one address; the proxy\index{proxy} routes each one by its model id to the backend that serves it, copies the answer back byte for byte, and records tokens, time to first token, decode speed, cached tokens and cost in one ledger. It needs Python 3.11 or newer and one package, \code{cryptography}. No GPU, no torch.

\section{Installing from a checkout}

\begin{shell}
git clone https://github.com/TandemLLM/TandemApp.git && cd TandemApp
pip install '.[host]'  # the host extra adds psutil and NVML for the System view
tandem-app                 # http://127.0.0.1:8080/dashboard/ opens on Setup
\end{shell}

The distribution is called \code{tandemapp}. Its \code{host} extra adds \code{psutil} and \code{nvidia-ml-py}, which the System view uses to read the machine; without them the collector falls back to \code{/proc} and \code{nvidia-smi}. The package is not on PyPI; install it from the checkout, or build the wheel with \code{python -m build --wheel .}. Its page (\code{dashboard/dist}) is a build output and not in git: build it first with \code{python3 tools/dashboard_bundle.py build} (Node 20.19 or newer), or fetch a released bundle with its \code{fetch} command and the bundle's sha256. A wheel build without the page stops and says how to get it. \code{python -m proxy} in the checkout does the same as \code{tandem-app} without installing anything but \code{cryptography}.

The page opens on Setup\index{Setup screen} while no backend exists: add one, test the connection, save. The panels that only an engine fills (caches, speculation, the engine side of the conductor) say that this server has no engine.

\section{The tandem-app command}\index{tandem-app}

\begin{table}[htbp]
\centering\small
\caption{The subcommands of \texttt{tandem-app}.}\label{tab:tandem-app}
\begin{tabularx}{\linewidth}{@{}L{.3\linewidth}Y@{}}
\toprule
\hd{command} & \hd{what it does} \\
\midrule
\code{tandem-app}, \code{tandem-app serve} & the server: the proxy, the dashboard and the ledger; every flag of Table~\ref{tab:ref-proxy-flags} \\
\code{tandem-app check} & reads the environment and a copy of the config database through the schema and prints what would run: backends, key hints, locked fields, routing rules, the listen address. Writes nothing; exits 1 on an error that names the variable \\
\code{tandem-app secrets init} & makes the master key file beside the config database (mode 0600) when no master key answers; \code{secrets status} says which source answers \\
\code{tandem-app backends list} & one line per backend, from a copy of the database \\
\code{tandem-app hostcollect} & the remote host collector (Section~\ref{sec:hostcollect}) \\
\code{tandem-app launch vllm|sglang} & starts vLLM or SGLang with TandemLLM's plugin and registers it as a backend (Section~\ref{sec:launch}) \\
\code{tandem-app version} & the installed version \\
\bottomrule
\end{tabularx}
\end{table}

\section{A first backend}

A backend can be added in the GUI or declared in an env file. The file form suits a headless machine or a container. A local vLLM:

\begin{shell}
cat > app.env <<'EOF'
TANDEM_BACKENDS=vllm1
TANDEM_BACKEND_VLLM1_KIND=vllm
TANDEM_BACKEND_VLLM1_URL=http://127.0.0.1:8100/v1
EOF
tandem-app check --env-file app.env        # prints what would run, writes nothing
tandem-app serve --env-file app.env
\end{shell}

A cloud API, with the key in the environment, where it stays in memory and is never written to the database:

\begin{shell}
printf 'TANDEM_BACKENDS=deepseek\nTANDEM_BACKEND_DEEPSEEK_KIND=deepseek\nTANDEM_BACKEND_DEEPSEEK_KEY=%s\n' \
    "$DEEPSEEK_API_KEY" > app.env && chmod 600 app.env
tandem-app serve --env-file app.env &
curl -s localhost:8080/v1/chat/completions -H 'Content-Type: application/json' \
    -d '{"model": "deepseek-flash", "messages": [{"role": "user", "content": "Say hi"}]}'
\end{shell}

Then open \code{http://127.0.0.1:8080/dashboard/#/usage}: the request is one row there, with its tokens, its time to first token and its cost at the model's price. Chapter~\ref{ch:other-engines} has a section for each kind of backend.

\section{In docker}\index{docker}

\code{packaging/Dockerfile.app} builds an image on \code{python:3.12-slim} from the same wheel, with the host extra and the built page inside, and no compiler or torch. It is configured by an env file; \code{packaging/app.env.example} is one to start from.

\begin{shell}
docker build -f packaging/Dockerfile.app -t tandemllm-app app
cp packaging/app.env.example app.env  # fill in; no quotes around values
docker run -d --name tandem-app -p 127.0.0.1:8080:8080 -v tandem-app:/data \
    --env-file app.env tandemllm-app
\end{shell}

The config database, the usage ledger and the key file live in the \code{/data} volume. Inside the container the server binds \code{0.0.0.0}, so what the outside reaches is what \code{-p} publishes: \code{127.0.0.1:8080} for this machine only. Keep it that way: no admin password is set by default, and this release has no sign-in screen (\enginemanual{Access and sign-in}). The image runs as an unprivileged user and answers its health check on \code{/health}. The release's tests build and install the wheel; they do not build the image. A volume from 0.15.0 or before keeps its old name; \code{docs/proxy.md} says how to mount or copy it.

Inside a container the host collector sees the container, not the machine. Run it with \code{--pid=host} (and \code{--gpus all} for NVML) for the machine's view, or run \code{tandem-app hostcollect}\index{tandem-app hostcollect@\texttt{tandem-app hostcollect}} on the machine itself.

\section{Where the proxy keeps its files}

\begin{table}[htbp]
\centering\small
\caption{The standalone proxy's files.}\label{tab:proxy-files}
\begin{tabularx}{\linewidth}{@{}L{.18\linewidth}YL{.36\linewidth}@{}}
\toprule
\hd{file} & \hd{what} & \hd{default} \\
\midrule
config database & backends, encrypted keys, models, prices, routing rules, users, the fleet's nodes and saved window layouts, an audit table & \code{~/.tandemllm/config.sqlite3} (\code{--config-db}, \code{TANDEM_CONFIG_DB}) \\
usage ledger & one row per request: tokens, times, cost; no text & \code{~/.tandemllm/ledger.sqlite3} (\code{--usage-ledger}; \code{off} for none) \\
fleet store & the fleet's metric history, event log, alerts and rule edits & \code{fleet.sqlite3} beside the config database (\code{--fleet-store}; \code{off} for memory only) \\
master key file & the key the backends' keys are encrypted with, when no other source answers & \code{~/.tandemllm/secrets.key}, mode 0600 \\
conductor documents & the harness's profiles, router, packages and journal & \code{~/.tandemllm/conductor} (\code{--conductor-config}; \code{off} for none) \\
\bottomrule
\end{tabularx}
\end{table}

\section{The proxy inside the engine's server}

The engine's own server can relay too. Start it with \code{--backends <config database>} (\code{QSE_BACKENDS}\index{qsebackends@\texttt{QSE\_BACKENDS}} in \code{serve.env}; empty means off). It then serves its own model as before and relays every request whose model belongs to a backend. A relayed request runs on its handler's own thread and never takes the engine's lock or a place in its queue, so a slow cloud answer cannot hold up the engine, and the engine's queue cannot hold up a cloud answer. With \code{--backends}\index{backends@\texttt{-\kern0pt-backends}} off the engine loads none of the proxy's code and serves exactly as without it.

\section{The fleet}\label{sec:fleet}\index{fleet}

TandemApp keeps the registry of the Sparks it manages and answers the fleet API, \code{/v1/fleet/*}, which the dashboard's shell, scope bar and windows read. The contract is \code{docs/contract/fleet-v1}. The browser talks to \code{tandem-app} only: every figure of a node is read on the server.

On one box the fleet has one node: the machine \code{tandem-app} runs on, when it has an NVIDIA GPU or an engine of its own. The first start registers it with its hostname as the alias, the first identity colour and an id from the machine (the DMI serial, else the machine id), never from its address; a new address keeps the node and is noted in the log. A machine in front of cloud APIs only, a Mac for instance, is the controller's host and no node. When TandemNode runs on the same machine (Section~\ref{sec:tandemnode}), the controller reads the machine through it, the path every other node takes; while it does not answer, in-process. A node on another machine is read through its agent once pairing has given the controller its trust, and is listed as off until then.

\begin{table}[htbp]
\centering\small
\caption{The fleet API on a fleet of one.}\label{tab:fleet-one}
\begin{tabularx}{\linewidth}{@{}L{.26\linewidth}Y@{}}
\toprule
\hd{route} & \hd{what it answers} \\
\midrule
\code{nodes}, \code{nodes/<id>} & the node: status, version and inventory (OS, kernel, driver, CUDA, ConnectX ports, disk, interfaces); \code{PATCH} sets the alias, the colour (1 to 8, one node each), the icon and the tags \\
\code{deployments} & each model an engine on the node serves itself, as a Single deployment, its status from the backend's health check \\
\code{links}, \code{topology} & the LAN of the default route, and the internet when a cloud backend exists \\
\code{metrics} & decode tok/s, requests and TTFT p50 from each engine (a TandemLLM engine's live snapshot, a vLLM's or SGLang's \code{/metrics}); GPU, power, temperature, clocks, memory, CPU, disk, LAN and fabric from the host collector, kept in the fleet store (Section~\ref{sec:fleet-store}). The hardware is sampled every 2\,s; a TandemLLM engine's own one-second samples are read every 2\,s while a dashboard reads the fleet and every 30\,s otherwise, the first read with the engine's last five minutes. A point without a sample is null \\
\code{memory} & the engine's split (weights, KV, caches) beside the rest of the machine's use and the page cache, GiB \\
\code{logs} & \code{tandem-app}'s own log and its TandemLLM engines', merged by time; info and up by default, the HTTP access lines only with \code{access=1} \\
\code{layouts/<page>} & a page's window layout per device class, in the config database; the last write wins \\
\code{events} & the event log, newest first, by time, node and type \\
\code{alerts}, \code{rules} & the alerts firing now and their acknowledgement; the alert rules and their edits \\
\bottomrule
\end{tabularx}
\end{table}

Reads follow the dashboard rule and writes the admin rule (\enginemanual{Access and sign-in}), as for \code{/v1/config}. The engine's own server, which serves the dashboard on a box without \code{tandem-app}, has no fleet API; the dashboard then shows the box alone.

\subsection{History, events and alert rules}\label{sec:fleet-store}

The fleet store is one SQLite file beside the config database, \code{fleet.sqlite3} (\code{--fleet-store}, \code{TANDEM_FLEET_STORE}; \code{off} keeps memory only). Each metric is kept in slots of 2\,s for 15 minutes (in memory, written at a clean stop), 30\,s for 26 hours, and 10 minutes for the retention, 30 days by default and at least 7 (\code{--fleet-retention-days}). The 5-minute range reads the 2\,s slots, 1 hour to 24 hours the 30\,s slots, and 7 days the 30\,s slots for its newest day and the 10-minute slots before. Changed slots are written every 5 minutes in one transaction; a crash loses at most those minutes, a clean stop nothing. TTFT is the median of the requests in each window, not a mean.

The event log records nodes going off, draining and up again, pairing, take-over and release, a new address, an engine down, up, started or restarted, a deployment's status, each alert firing, resolving and acknowledged, and the controller's start and stop.

\begin{table}[htbp]
\centering\small
\caption{The alert rules a fleet starts with, set for a GB10 (unified memory, no fan or power cap exposed).}\label{tab:fleet-rules}
\begin{tabularx}{\linewidth}{@{}L{.18\linewidth}YL{.1\linewidth}L{.24\linewidth}@{}}
\toprule
\hd{rule} & \hd{fires} & \hd{severity} & \hd{resolves} \\
\midrule
\code{gpu-temp} & GPU at 85\,\textdegree{}C or more for 60\,s & err & below 80\,\textdegree{}C for 60\,s \\
\code{throttle} & a thermal or power-brake clock event for 30\,s & warn & none for 60\,s \\
\code{mem-avail} & MemAvailable under 10\,GiB for 30\,s & warn & above 12\,GiB for 30\,s \\
\code{disk-free} & a mount under 10\,\% free for 5 minutes & warn & above 12\,\% for 5 minutes \\
\code{node-down} & no answer for 30\,s & err & up for 10\,s \\
\code{engine-down} & an engine's health check fails for 60\,s & err & one check that passes \\
\code{host-sleep} & the controller's machine may sleep (a Mac) & warn & \\
\bottomrule
\end{tabularx}
\end{table}

\code{PATCH /v1/fleet/rules/<id>} turns a rule on or off and sets its notify targets, threshold, hold and scope (all nodes or one). An alert names the page to open and stays listed, acknowledged or not, until it resolves. Notify \code{app} is the alerts list and \code{system} is for the page to show; the controller makes no outbound call.

\section{TandemNode}\label{sec:tandemnode}\index{TandemNode}

TandemNode is the agent on each Spark that TandemApp's controller talks to. It ships in the same package as \code{tandem-node}; on an \code{install.sh} install \code{tandem node} runs it with the install's engine as its engine. It listens on port 7700.

\begin{table}[htbp]
\centering\small
\caption{TandemNode's routes.}\label{tab:tandemnode}
\begin{tabularx}{\linewidth}{@{}L{.3\linewidth}Y@{}}
\toprule
\hd{route} & \hd{what it answers} \\
\midrule
\code{GET /node/identity}, \code{/node/inventory} & node id, DMI product and serial, OS, kernel, driver, CUDA, TandemLLM version; disks, interfaces (ConnectX ports with firmware and RoCE state), memory, cores, GPUs \\
\code{GET /node/metrics} & an event a second on the whole second: GPU, unified memory split, per-core CPU, disks, LAN and fabric per interface, thermal zones, fans. A field the board does not expose is absent \\
\code{/node/engines} & the engines, their health and jobs; \code{POST .../start}, \code{stop}, \code{restart} for an engine the agent manages; the engine's own read routes and live stream relayed \\
\code{/node/weights} & the Setup Guide's weights with a content id per pack; \code{DELETE} removes one as the Setup Guide does \\
\code{GET /node/logs} & the agent's and its engines' lines, by level, text and source \\
\code{POST /node/linktest} & to another node: RTT, bandwidth per interface, MTU, RoCE state, a port running at half its rate \\
\code{/node/maintenance} & whether the machine's reboot and update are set up here (simulated on a virtual node); \code{POST .../reboot}, \code{update} run them \\
\code{/node/setup/*} & the engine's Setup Guide routes, unchanged \\
\bottomrule
\end{tabularx}
\end{table}

Every route needs its caller checked. Until pairing has filled the trust store (\code{~/.tandemllm/node/trust}), the agent listens on loopback only, in plain HTTP, and every request carries the loopback secret from \code{~/.tandemllm/node/secret} (mode 0600, made at the first start). Any other address needs mutual TLS: a client without a certificate the store vouches for fails the handshake. It never listens on \code{0.0.0.0} or a WAN address; a remote node is reached through an SSH tunnel to its loopback. \code{tandem-node check} says how it would listen. \code{tandem-node unit} prints a systemd user unit for the same options.

The machine's reboot and update are commands the node's owner sets in \code{~/.tandemllm/node/node.json}; neither is set by default, because a reboot needs a right the agent never takes by itself. A virtual node (\code{--virtual}, a second agent on one machine for tests) simulates both and never touches the machine.

\subsection{Maintenance: drain}\index{TandemNode!drain}

Drain on a Spark's page stops new requests to it through TandemApp's API: a model's route goes to its fallback, and with none the client gets a 503 with \code{Retry-After}, so it tries again later. Requests in flight finish; the Spark shows draining (a half-filled circle), then drained once nothing runs, and can be updated or rebooted. Undrain takes requests again. Every action that stops or removes something (restart the engine, reboot, update, forget, release a rented Spark) first says what goes and what stays, and runs only from that page's confirmation.

\subsection{Pairing}\index{TandemNode!pairing}

A node and a controller trust each other through pinned certificates; no password is stored. \code{tandem node pair} on the node prints a 6-digit code, valid 5 minutes for 5 tries and once; Add Spark in TandemApp takes the node's address and the code. The two run SPAKE2 (RFC 9382's flow in RFC 3526's 2048-bit group) with both certificates' fingerprints bound to the code, so a wrong code or a certificate swapped on the way pins nothing. Afterwards the node listens on that address in mutual TLS, owned by that one controller. SSH once is the other way in: the controller logs in a single time with a password or a key held in memory for that call, pins its certificate on the node and the node's on its side, and remembers the host key (a changed one is refused). Another TandemApp can take a node over only when the owner releases it, or with \code{tandem node takeover} on the node, or with SSH access. \code{tandem node reset} forgets the owner and every pin. Pairing and release need the admin password, or a caller on the controller's machine.

\section{Access}

The standalone proxy follows the engine's access rule (\enginemanual{Access and sign-in}) with its own variable, \code{TANDEM_ADMIN_TOKEN}, and adds user keys for the OpenAI routes (Section~\ref{sec:users}). It binds \code{127.0.0.1} by default.
